\documentclass[11pt]{article}
\newcommand{\doctitle}{Assessing Carryover in the Analysis Model:
Standard Methodology and Application to pmsimstats}
\newcommand{\shorttitle}{Carryover Analysis Assessment}
\newcommand{\docdate}{2026-03-22}
% pmsimstats-ng standard document preamble
% Usage: % pmsimstats-ng standard document preamble
% Usage: % pmsimstats-ng standard document preamble
% Usage: \input{pmsimstats-preamble} after \documentclass[11pt]{article}
%
% Documents must define before \input:
%   \newcommand{\doctitle}{...}       Full document title
%   \newcommand{\shorttitle}{...}     Short title for header
%   \newcommand{\docdate}{YYYY-MM-DD} Document date

\usepackage[margin=1in]{geometry}
\usepackage{amsmath,amssymb}
\usepackage[T1]{fontenc}
\usepackage{lmodern}
\usepackage{microtype}
\usepackage{graphicx}
\graphicspath{{figures/}}
\usepackage{booktabs}
\usepackage{longtable}
\usepackage{array}
\usepackage{hyperref}
\usepackage{xcolor}
\usepackage{fancyhdr}
\usepackage{parskip}
\usepackage{caption}
\usepackage{enumitem}
\usepackage{verbatim}
\usepackage{listings}

\lstset{
  language=R,
  basicstyle=\small\ttfamily,
  keywordstyle=\color{blue!70!black},
  commentstyle=\color{green!50!black},
  stringstyle=\color{red!60!black},
  breaklines=true,
  frame=single,
  framerule=0.5pt,
  rulecolor=\color{gray!40},
  backgroundcolor=\color{gray!5},
  xleftmargin=1em,
  framexleftmargin=1em,
  columns=flexible
}

\hypersetup{
  colorlinks=true,
  linkcolor=blue!60!black,
  citecolor=green!50!black,
  urlcolor=blue!60!black
}

\pagestyle{fancy}
\fancyhf{}
\lhead{\footnotesize pmsimstats-ng Technical Report}
\rhead{\footnotesize \shorttitle}
\rfoot{\footnotesize \thepage}
\renewcommand{\headrulewidth}{0.4pt}

\title{\doctitle}
\author{pmsimstats team}
\date{\docdate}
 after \documentclass[11pt]{article}
%
% Documents must define before \input:
%   \newcommand{\doctitle}{...}       Full document title
%   \newcommand{\shorttitle}{...}     Short title for header
%   \newcommand{\docdate}{YYYY-MM-DD} Document date

\usepackage[margin=1in]{geometry}
\usepackage{amsmath,amssymb}
\usepackage[T1]{fontenc}
\usepackage{lmodern}
\usepackage{microtype}
\usepackage{graphicx}
\graphicspath{{figures/}}
\usepackage{booktabs}
\usepackage{longtable}
\usepackage{array}
\usepackage{hyperref}
\usepackage{xcolor}
\usepackage{fancyhdr}
\usepackage{parskip}
\usepackage{caption}
\usepackage{enumitem}
\usepackage{verbatim}
\usepackage{listings}

\lstset{
  language=R,
  basicstyle=\small\ttfamily,
  keywordstyle=\color{blue!70!black},
  commentstyle=\color{green!50!black},
  stringstyle=\color{red!60!black},
  breaklines=true,
  frame=single,
  framerule=0.5pt,
  rulecolor=\color{gray!40},
  backgroundcolor=\color{gray!5},
  xleftmargin=1em,
  framexleftmargin=1em,
  columns=flexible
}

\hypersetup{
  colorlinks=true,
  linkcolor=blue!60!black,
  citecolor=green!50!black,
  urlcolor=blue!60!black
}

\pagestyle{fancy}
\fancyhf{}
\lhead{\footnotesize pmsimstats-ng Technical Report}
\rhead{\footnotesize \shorttitle}
\rfoot{\footnotesize \thepage}
\renewcommand{\headrulewidth}{0.4pt}

\title{\doctitle}
\author{pmsimstats team}
\date{\docdate}
 after \documentclass[11pt]{article}
%
% Documents must define before \input:
%   \newcommand{\doctitle}{...}       Full document title
%   \newcommand{\shorttitle}{...}     Short title for header
%   \newcommand{\docdate}{YYYY-MM-DD} Document date

\usepackage[margin=1in]{geometry}
\usepackage{amsmath,amssymb}
\usepackage[T1]{fontenc}
\usepackage{lmodern}
\usepackage{microtype}
\usepackage{graphicx}
\graphicspath{{figures/}}
\usepackage{booktabs}
\usepackage{longtable}
\usepackage{array}
\usepackage{hyperref}
\usepackage{xcolor}
\usepackage{fancyhdr}
\usepackage{parskip}
\usepackage{caption}
\usepackage{enumitem}
\usepackage{verbatim}
\usepackage{listings}

\lstset{
  language=R,
  basicstyle=\small\ttfamily,
  keywordstyle=\color{blue!70!black},
  commentstyle=\color{green!50!black},
  stringstyle=\color{red!60!black},
  breaklines=true,
  frame=single,
  framerule=0.5pt,
  rulecolor=\color{gray!40},
  backgroundcolor=\color{gray!5},
  xleftmargin=1em,
  framexleftmargin=1em,
  columns=flexible
}

\hypersetup{
  colorlinks=true,
  linkcolor=blue!60!black,
  citecolor=green!50!black,
  urlcolor=blue!60!black
}

\pagestyle{fancy}
\fancyhf{}
\lhead{\footnotesize pmsimstats-ng Technical Report}
\rhead{\footnotesize \shorttitle}
\rfoot{\footnotesize \thepage}
\renewcommand{\headrulewidth}{0.4pt}

\title{\doctitle}
\author{pmsimstats team}
\date{\docdate}


\begin{document}
\maketitle
\thispagestyle{fancy}
\tableofcontents

\textbf{pmsimstats team}

\bigskip\noindent\rule{\textwidth}{0.4pt}\bigskip

\section{The Question}

When simulating clinical trial data with carryover effects in the
data-generating process (DGP), what is the proper way to handle
carryover in the analysis model? The pmsimstats framework currently
offers two options -- a continuous drug indicator (\texttt{Dbc}) that
bakes carryover into the treatment variable, and a separate \texttt{tsd}
covariate -- but the choice between them (or whether to include
carryover at all) has not been systematically evaluated.

\bigskip\noindent\rule{\textwidth}{0.4pt}\bigskip

\section{The 2x2 Factorial Framework}

The standard simulation framework treats the DGP and analysis
model as a factorial design with two factors:

\begin{table}[h]
\centering
\begin{tabular}{lll}
\toprule
 & Analysis ignores carryover & Analysis models carryover \\
\midrule
\textbf{DGP: no carryover} & Baseline (Type I error check) & Cost of unnecessary adjustment \\
\textbf{DGP: carryover present} & Cost of ignoring carryover & Correct specification \\
\bottomrule
\end{tabular}
\end{table}

Data are generated under each DGP condition, analyzed under each
analysis model, and the following metrics are reported for all
four cells:

\begin{itemize}
  \item Type I error (rejection rate when $c_{bm} = 0$)
  \item Power (rejection rate when $c_{bm} > 0$)
  \item Bias of the treatment effect estimate
  \item Coverage of the 95\% confidence interval
\end{itemize}

The diagonal cells (matched DGP and analysis) should perform
well. The off-diagonal cells reveal the consequences of model
misspecification: inflated Type I error when carryover is ignored,
or unnecessary power loss when carryover is modeled but absent.

\bigskip\noindent\rule{\textwidth}{0.4pt}\bigskip

\section{What the Literature Says}

\subsection{The definitive reference}

Senn (2002), \textit{Cross-over Trials in Clinical Research}, 2nd
edition, provides the authoritative treatment. The key conclusions
for simulation design:

\begin{enumerate}
  \item \textbf{The two-stage procedure is invalid.} The older approach --
    test for carryover first, then decide whether to include it
    in the treatment model -- was shown to be statistically
    unsound. The carryover test has low power (crossover designs
    are not designed to detect carryover; they are designed to
    detect treatment effects). Conditioning the treatment analysis
    on the outcome of the carryover test inflates the Type I error
    of the treatment test beyond the nominal level.

  \item \textbf{The modern consensus:} If carryover is scientifically
    plausible, either (a) design the trial to eliminate it via a
    washout period of sufficient length, or (b) always include
    carryover in the model and accept the associated power loss.
    Do not test-then-decide.

  \item \textbf{For simulation studies:} Generate data under a range of
    carryover magnitudes (including zero), analyze under models
    that include and exclude the carryover term, and report
    operating characteristics for each combination.
\end{enumerate}

\subsection{Supporting references}

\begin{itemize}
  \item Jones \& Kenward (2014), \textit{Design and Analysis of Cross-Over
    Trials}, 3rd edition. Extends Senn's framework to higher-order
    crossover designs and provides simulation templates for
    evaluating carryover model specifications.

  \item Zucker et al. (1997), `Combining single patient (N-of-1)
    trials to estimate population treatment effects and its
    application to crossover designs.' Establishes the mixed-effects
    framework for aggregated N-of-1 trials, which is the basis of
    the pmsimstats analysis model.

  \item Araujo et al. (2016), `Understanding variation in sets of
    N-of-1 trials.' Addresses heterogeneity in treatment effects
    across participants, which is precisely the biomarker-treatment
    interaction that pmsimstats is designed to detect.
\end{itemize}

\bigskip\noindent\rule{\textwidth}{0.4pt}\bigskip

\section{How This Maps to pmsimstats}

\subsection{Carryover in the DGP}

The pmsimstats DGP introduces carryover through the mean
structure of the biological response (BR) component. During
off-drug periods, the BR mean at timepoint $t$ is augmented by
the decayed effect from the previous timepoint:

\[
\mu_{BR}(t) = \text{base}(t) + \mu_{BR}(t-1) \cdot
(1/2)^{t_{sd} / t_{1/2}}
\]

where $t_{sd}$ is the time since drug discontinuation and
$t_{1/2}$ is the carryover half-life. This is an exponential
decay model: the residual drug effect halves every $t_{1/2}$
weeks.

Additionally, the BM-BR correlation decays during off-drug
periods at a rate tied to the same half-life:

\[
\text{Cor}(bm, br_t) = c_{bm} \cdot e^{-\lambda_{cor} \cdot
t_{sd}}, \quad \lambda_{cor} = \ln(2) / t_{1/2}
\]

\subsection{Carryover in the analysis model}

The current codebase offers two approaches:

\textbf{Approach A: Continuous drug indicator (\texttt{Dbc}).} The binary
on/off-drug indicator is replaced with a continuous version:

\begin{itemize}
  \item On drug: $Dbc = 1$
  \item Off drug: $Dbc = (1/2)^{s \cdot t_{sd} / t_{1/2}}$
\end{itemize}

The model formula becomes:
\texttt{Sx \~{} bm + t + Dbc + bm:Dbc}, with the interaction \texttt{bm:Dbc} as
the test statistic.

This approach conflates treatment status and carryover into a
single variable. The \texttt{bm:Dbc} coefficient captures both the
treatment-biomarker interaction and any residual
carryover-biomarker interaction. Whether this is a feature or a
bug depends on the scientific question: if the goal is to detect
`does the biomarker predict who benefits from drug exposure
(including residual exposure),' then \texttt{Dbc} is appropriate. If
the goal is to detect `does the biomarker predict the on-drug
treatment effect, adjusted for carryover,' then a separate
carryover term is needed.

\textbf{Approach B: Separate carryover covariate (\texttt{tsd}).} A linear
time-since-discontinuation term is added to the model:

\texttt{Sx \~{} bm + t + Db + bm:Db + tsd}

where \texttt{Db} is binary (on/off drug) and \texttt{tsd} absorbs the
carryover effect. The interaction \texttt{bm:Db} then estimates the
biomarker-treatment interaction after adjusting for carryover.
This is closer to the standard crossover trial approach but uses
a linear carryover effect, which is a mild misspecification when
the DGP uses exponential decay.

\textbf{Approach C: Ignore carryover.} The model uses only the binary
drug indicator:

\texttt{Sx \~{} bm + t + Db + bm:Db}

Off-drug observations with residual drug effect are treated as
if no drug were present. This biases the treatment contrast
toward zero (carryover makes off-drug periods look more like
on-drug periods), reducing power but potentially inflating
Type I error if the bias is systematic.

\subsection{The non-standard element}

The \texttt{Dbc} parameterization (Approach A) is creative but
non-standard in the crossover trial literature. Standard practice
separates the treatment indicator from the carryover adjustment.
The \texttt{Dbc} approach would require explicit justification in a
methods paper, ideally supported by the factorial simulation
comparison described below.

\bigskip\noindent\rule{\textwidth}{0.4pt}\bigskip

\section{Proposed Simulation Study}

A rigorous evaluation would compare three analysis models across
the carryover DGP sweep.

\subsection{Design}

\textbf{DGP factor:} $t_{1/2} \in \{0, 0.25, 0.5, 1.0, 2.0\}$ weeks

\textbf{Analysis factor:} Three models:

\begin{enumerate}
  \item \textbf{Ignore:} \texttt{Sx \~{} bm + t + Db + bm:Db}
    (binary drug indicator, no carryover term)
  \item \textbf{Separate:} \texttt{Sx \~{} bm + t + Db + bm:Db + tsd}
    (binary drug indicator + linear carryover covariate)
  \item \textbf{Continuous:} \texttt{Sx \~{} bm + t + Dbc + bm:Dbc}
    (exponential-decay drug indicator, current default)
\end{enumerate}

\textbf{Fixed parameters:} $N = 70$, $\rho = 0.7$, all four trial
designs.

\textbf{Biomarker correlation:} $c_{bm} \in \{0, 0.3, 0.45\}$
(null, moderate, strong).

\textbf{Replicates:} 1,000 per cell.

\subsection{Metrics}

For each of the $5 \times 3 \times 3 \times 4 = 180$ cells:

\begin{itemize}
  \item \textbf{Type I error} (at $c_{bm} = 0$): should be $\leq 0.05$
  \item \textbf{Power} (at $c_{bm} > 0$): higher is better
  \item \textbf{Bias:} $E[\hat{\beta}] - \beta_{true}$
  \item \textbf{RMSE:} $\sqrt{E[(\hat{\beta} - \beta_{true})^2]}$
  \item \textbf{Coverage:} proportion of 95\% CIs containing $\beta_{true}$
\end{itemize}

\subsection{Expected findings}

Based on the literature and the DGP structure:

\begin{itemize}
  \item \textbf{Ignore (model 1)} will have correct Type I error at
    $t_{1/2} = 0$ but may have inflated Type I error at higher
    carryover values (residual drug effect creates a systematic
    correlation between biomarker and off-drug response that the
    model attributes to noise). Power will decline with increasing
    carryover.

  \item \textbf{Separate (model 2)} will have correct Type I error across
    all carryover values (the \texttt{tsd} term absorbs the carryover
    signal). Power will be slightly lower than model 1 at
    $t_{1/2} = 0$ (cost of the extra parameter) but higher at
    $t_{1/2} > 0$ (correct adjustment recovers treatment
    contrast). Mild bias possible due to linear vs exponential
    misspecification.

  \item \textbf{Continuous (model 3)} will have correct Type I error if the
    \texttt{Dbc} decay rate matches the DGP. Power behavior depends on
    whether the conflation of treatment and carryover helps or
    hurts: it may outperform model 2 when the carryover half-life
    is correctly specified, but underperform when misspecified.
\end{itemize}

\subsection{Infrastructure}

The existing codebase supports all three models:

\begin{itemize}
  \item Model 1: \texttt{op\$carryover\_t1half = 0, op\$simplecarryover = FALSE}
  \item Model 2: \texttt{op\$carryover\_t1half = 0, op\$simplecarryover = TRUE}
  \item Model 3: \texttt{op\$carryover\_t1half = [matching DGP value]}
\end{itemize}

The simulation grid would need to cross the DGP carryover
parameter with the analysis model choice. This requires a
modification to \texttt{generateSimulatedResults()} (or the 2025
equivalent) to accept multiple analysis parameter sets per
DGP configuration, which is already supported via the
\texttt{analysisparams} argument.

\bigskip\noindent\rule{\textwidth}{0.4pt}\bigskip

\section{Recommendation}

The factorial comparison described in Section 5 should be run
before finalizing the manuscript. It addresses two questions that
reviewers will ask:

\begin{enumerate}
  \item \textbf{Does the \texttt{Dbc} parameterization control Type I error?} If
    it does not, the power estimates reported in the paper are
    unreliable.

  \item \textbf{Is the power gain from modeling carryover worth the
    complexity?} If the `ignore' model performs nearly as well as
    the carryover-aware models, the simpler approach is
    preferable.
\end{enumerate}

The simulation infrastructure already exists. The primary effort
is configuring the parameter grid and collecting the results into
a comparison table. Estimated effort: 4--6 hours of computation
(parallelizable) plus 2--3 hours of analysis and writing.

\vfill
\noindent\rule{\textwidth}{0.4pt}
{\footnotesize
Rendered on 2026-03-25 at 19:00 PDT.\\
Source: \textasciitilde/prj/alz/10-pmsimstats-ng/pmsimstats-ng/docs/10-carryover-analysis-model-assessment.tex}
\end{document}
